% =====================================================================
%  Referencias, índices y bibliografía — suite de documentos LaTeX
%  Nivel 2 (Capacidades) · clase article.
%  Compilar: latexmk -pdf referencias-biblio.tex  (corre biber solo)
% =====================================================================
\documentclass[11pt,a4paper]{article}

% --- Base estándar ---
\usepackage[utf8]{inputenc}
\usepackage[T1]{fontenc}
\usepackage{lmodern}
\usepackage{microtype}
\usepackage[spanish,es-noshorthands]{babel}
\usepackage{csquotes}     % requerido por biblatex para comillas correctas
\usepackage{amsmath}\usepackage{amssymb}
\usepackage{booktabs}
\usepackage[margin=2.4cm]{geometry}
\usepackage{enumitem}
\usepackage{xcolor}
\usepackage{listings}
% --- Bibliografía moderna ---
\usepackage[backend=biber, style=numeric, sorting=nyt]{biblatex}
\addbibresource{refs.bib}
% --- Enlaces y referencias cruzadas ---
\usepackage{hyperref}
\hypersetup{unicode, pdftitle={Referencias, índices y bibliografía},
            pdfauthor={Rodrigo Martín Vidal Galán},
            colorlinks=true, linkcolor=blue!50!black,
            citecolor=green!45!black, urlcolor=blue!50!black}
\usepackage[spanish,capitalise,nameinlink]{cleveref}

% --- Estilo de código ---
\definecolor{codeback}{RGB}{245,247,250}
\definecolor{codekw}{RGB}{20,60,120}
\definecolor{codecom}{RGB}{110,120,135}
\lstset{
  basicstyle=\ttfamily\footnotesize,
  keywordstyle=\color{codekw}\bfseries,
  commentstyle=\color{codecom}\itshape,
  backgroundcolor=\color{codeback},
  frame=single, rulecolor=\color{codekw!30},
  breaklines=true, showstringspaces=false, columns=fullflexible,
  xleftmargin=8pt, framesep=4pt, extendedchars=true,
  literate={á}{{\'a}}1 {é}{{\'e}}1 {í}{{\'i}}1 {ó}{{\'o}}1 {ú}{{\'u}}1
           {ñ}{{\~n}}1 {¿}{{?`}}1 {¡}{{!`}}1,
}
\newcommand{\C}[1]{\texttt{\textbackslash #1}}
\newcommand{\pkg}[1]{\textsf{#1}}

\title{\textbf{Referencias, índices y bibliografía}\\[2pt]
       \large El \enquote{aparato} del documento}
\author{Rodrigo Martín Vidal Galán}
\date{Junio de 2026}

\begin{document}
\maketitle

\noindent Todo lo que \textbf{conecta} el documento: referencias cruzadas
internas, enlaces, notas, índices y la \textbf{bibliografía}. La idea común es la
misma que en toda la suite: \textbf{nunca escribir un número a mano}; dejá que
\LaTeX{} lo resuelva y se renumere solo. \emph{(El índice analítico se trata en
\emph{Documentos largos}.)}

\newpage
\tableofcontents
\newpage

% =====================================================================
\section{Referencias cruzadas internas}
% =====================================================================
El mecanismo base son tres comandos: \C{label\{clave\}} marca un lugar (sección,
ecuación, figura\dots) y \C{ref\{clave\}}/\C{pageref\{clave\}} lo citan por
número o página:
\begin{lstlisting}
\section{Resultados}\label{sec:res}
...como vimos en la seccion~\ref{sec:res} (pagina~\pageref{sec:res}).
\end{lstlisting}
\begin{quote}\small
\textbf{Importante:} el \C{label} va \emph{después} del \C{caption} (en figuras/
tablas) o pegado al título de sección. Y hay que \textbf{compilar dos veces}
(o usar \texttt{latexmk}) para que los números se resuelvan; si ves \enquote{??}
en el PDF, falta una pasada.
\end{quote}

\paragraph{Mejor: \texttt{cleveref}.} En vez de escribir \enquote{la
figura}~\C{ref\{\dots\}} a mano, \C{cref} pone \textbf{el tipo y el número}
juntos, en el idioma correcto:
\begin{lstlisting}
\usepackage[spanish,capitalise,nameinlink]{cleveref} % DESPUES de hyperref
...
\cref{fig:a}        % -> figura 3
\Cref{fig:a}        % -> Figura 3  (al inicio de oracion)
\cref{eq:1,eq:2}    % -> ecuaciones 1 y 2  (agrupa solo)
\end{lstlisting}
\begin{center}\small
\renewcommand{\arraystretch}{1.25}
\begin{tabular}{@{}lll@{}}
\toprule
\textbf{Comando} & \textbf{Da} & \textbf{Cuándo} \\
\midrule
\C{ref}      & \texttt{3}            & número pelado \\
\C{eqref}    & \texttt{(3)}          & ecuaciones (con paréntesis) \\
\C{cref}     & \texttt{figura 3}     & \textbf{recomendado} (tipo + número) \\
\C{autoref}  & \texttt{figura 3}     & similar (de \pkg{hyperref}) \\
\C{pageref}  & \texttt{7}            & número de página \\
\bottomrule
\end{tabular}
\end{center}

\paragraph{Convención de claves.} Prefijá la clave según el tipo: \texttt{sec:},
\texttt{fig:}, \texttt{tab:}, \texttt{eq:}, \texttt{cap:}, \texttt{thm:}. Hace el
fuente legible y evita choques.

% =====================================================================
\section{Enlaces, marcadores y metadatos: \texttt{hyperref}}
% =====================================================================
\pkg{hyperref} convierte todas las referencias en \textbf{enlaces clicables},
genera los \textbf{marcadores} del PDF y guarda los metadatos. Se carga
\textbf{casi al final} (única excepción: \pkg{cleveref}, que va después):
\begin{lstlisting}
\usepackage{hyperref}
\hypersetup{
  unicode,                       % marcadores con acentos correctos
  pdftitle={Mi documento}, pdfauthor={Mi Nombre},
  colorlinks=true, linkcolor=blue, citecolor=green, urlcolor=blue,
  % hidelinks                    % alternativa: enlaces sin color (para imprimir)
}
\end{lstlisting}
URLs y enlaces externos:
\begin{lstlisting}
\url{https://ctan.org}                 % muestra la URL
\href{https://ctan.org}{el repositorio CTAN}  % texto con enlace
\end{lstlisting}
\begin{quote}\small
\textbf{Siempre} \C{hypersetup\{unicode\}}: sin eso, los marcadores del PDF
muestran símbolos raros en acentos y \enquote{¿}.
\end{quote}

% =====================================================================
\section{Notas e índices de contenido}
% =====================================================================
\begin{itemize}[leftmargin=*]
  \item \textbf{Notas al pie:} \C{footnote\{texto\}}.\footnote{Así se ve una nota
        al pie.} Se numeran solas.
  \item \textbf{Índice general:} \C{tableofcontents}; \C{listoffigures} y
        \C{listoftables} para los índices de figuras/tablas.
  \item \textbf{Agregar algo sin numerar al índice} (una sección \texttt{*}, la
        bibliografía): \C{addcontentsline\{toc\}\{section\}\{Bibliografía\}}.
  \item \textbf{Profundidad:} \C{setcounter\{tocdepth\}\{2\}} controla hasta qué
        nivel baja el índice.
  \item \textbf{Glosario / siglas:} paquete \pkg{glossaries} (\C{gls},
        \C{newacronym}); el \textbf{índice analítico}, \pkg{makeidx} (ver
        \emph{Documentos largos}).
\end{itemize}

% =====================================================================
\section{Bibliografía: el enfoque moderno (\texttt{biblatex} + biber)}
% =====================================================================
La forma recomendada hoy es \pkg{biblatex} con el motor \textbf{biber}: flexible,
internacionalizado y con muchos estilos. Tres pasos:

\paragraph{1) Un archivo \texttt{.bib}} con las fuentes (lo exportan Zotero,
Mendeley o JabRef):
\begin{lstlisting}
@book{spivak,
  author = {Spivak, Michael}, title = {Calculus},
  edition = {4}, publisher = {Publish or Perish}, year = {2008},
}
@article{shannon1948,
  author = {Shannon, Claude E.},
  title = {A Mathematical Theory of Communication},
  journal = {The Bell System Technical Journal},
  volume = {27}, pages = {379--423}, year = {1948},
}
\end{lstlisting}
Cada entrada tiene un \textbf{tipo} (\texttt{@book}, \texttt{@article},
\texttt{@online}\dots), una \textbf{clave} (\texttt{spivak}) y \textbf{campos}.

\paragraph{2) En el preámbulo:}
\begin{lstlisting}
\usepackage[backend=biber, style=numeric, sorting=nyt]{biblatex}
\addbibresource{refs.bib}
\end{lstlisting}

\paragraph{3) Citar y listar.} Se cita en el texto y, al final, \C{printbibliography}
arma la lista solo con lo citado:
\begin{center}\small
\renewcommand{\arraystretch}{1.25}
\begin{tabular}{@{}ll@{}}
\toprule
\textbf{Comando} & \textbf{Produce (estilo numeric)} \\
\midrule
\C{cite\{spivak\}}      & \cite{spivak} \\
\C{parencite\{spivak\}} & \parencite{spivak} \\
\C{textcite\{shannon1948\}} & \textcite{shannon1948} \\
\C{footcite\{ctan\}}    & nota al pie con la cita \\
\bottomrule
\end{tabular}
\end{center}
\noindent Ejemplo real en una oración: el cálculo se estudia en \textcite{spivak},
y la teoría de la información nace con \textcite{shannon1948}; el repositorio
central es CTAN \parencite{ctan}.

\paragraph{Estilos.} La opción \texttt{style=} cambia todo el formato sin tocar el
texto: \texttt{numeric} ([1], [2]), \texttt{alphabetic} ([Spi08]),
\texttt{authoryear} (Spivak 2008), \texttt{ieee}, \texttt{apa}\dots Cambiás una
palabra y se reformatea la bibliografía entera.

% =====================================================================
\section{El enfoque clásico (BibTeX + \texttt{natbib})}
% =====================================================================
El método antiguo usa \textbf{BibTeX} (no biber) y, para ciencias, el paquete
\pkg{natbib}:
\begin{lstlisting}
\usepackage[numbers]{natbib}
...
\citep{spivak}   % (cita entre parentesis)
\citet{spivak}   % cita textual: Autor (anio)
...
\bibliographystyle{plainnat}
\bibliography{refs}   % sin la extension .bib
\end{lstlisting}
Sigue vivo porque \textbf{muchas revistas} traen plantillas con su \texttt{.bst}.
Para trabajo nuevo y propio, preferí \pkg{biblatex}; si una revista exige BibTeX,
usá el que pidan.

\begin{quote}\small
\textbf{Compilación:} con \pkg{biblatex} la secuencia es
\texttt{pdflatex → biber → pdflatex → pdflatex}; con BibTeX,
\texttt{pdflatex → bibtex → pdflatex → pdflatex}. En ambos casos,
\texttt{latexmk -pdf} lo hace solo.
\end{quote}

% =====================================================================
\section{Buenas prácticas}
% =====================================================================
\begin{itemize}[noitemsep,leftmargin=*]
  \item \C{label}/\C{cref} para \textbf{todo} lo numerado; nunca el número a mano.
  \item \pkg{biblatex}\,+\,biber para trabajo nuevo; un gestor (Zotero/JabRef)
        para mantener el \texttt{.bib}.
  \item \pkg{hyperref} casi al final; \pkg{cleveref} después; \texttt{unicode} en
        \C{hypersetup}.
  \item Claves de cita y de \C{label} \textbf{con prefijo} y estables.
  \item Compilá con \texttt{latexmk} para no olvidar las pasadas de biber.
\end{itemize}

% =====================================================================
\section*{Ver también}
% =====================================================================
\addcontentsline{toc}{section}{Ver también}
\begin{itemize}[noitemsep,leftmargin=*]
  \item \emph{Documentos largos} — índice analítico, índices de figuras/tablas.
  \item \emph{Estándares} — \pkg{hyperref}/\pkg{cleveref} en el preámbulo.
  \item Manuales: \texttt{texdoc biblatex}, \texttt{texdoc hyperref},
        \texttt{texdoc cleveref}.
\end{itemize}

\printbibliography[heading=bibintoc, title={Bibliografía}]

\end{document}
